% Section 4: The bridge between forecast loss and decision loss
\subsection{The bridge between forecast loss and decision loss}
\label{sec:bridge}

The arguments of Sections~\ref{sec:rl_for_forecasting} and~\ref{sec:forecasting_for_rl}
converge on a single methodological question, namely what loss should govern the
training of the forecaster when the forecast is destined for use inside a decision.
The classical answer was given by \citet{GrangerPesaran2000}. The forecast
evaluation problem and the decision evaluation problem are not separable in general.
For a two-state two-action setup with payoff matrix $\pi$ and forecast $\hat p$
for the binary event $y$, the optimal action depends on whether $\hat p$ exceeds a
threshold determined by the relative payoffs. The mean squared error of the
forecast is monotone in $\lvert \hat p - p \rvert$ uniformly, but the realised
payoff is sensitive to errors only in the neighbourhood of the threshold. Two
forecasts with identical MSE rankings can swap under decision-theoretic ranking
once the cost structure is asymmetric, and \citet{GrangerPesaran2000} construct
the explicit reversal example. The bridge between meteorological forecast scoring
via the Kuipers score and financial market timing tests is one of their key
contributions.

The statistical complement is \citet{GneitingRaftery2007}. A scoring rule
$S(F, y)$ is proper if the forecaster's expected score is maximised at the true
forecast distribution $F = G$, and strictly proper if the maximiser is unique. They
show that the log score, the continuous ranked probability score, and the energy
score are strictly proper, that the squared error loss is not strictly proper in
distributional terms but is consistent for the mean, and that every strictly proper
rule arises from a Bregman divergence under a convex score-generating function.
Strictly proper rules are the right starting point for distribution forecasting
because they prevent the forecaster from gaming the score by hedging.

The decision-focused programme of \citet{ElmachtoubGrigas2022} and
\citet{DontiAmosKolter2017} can be read as the Wald (1950) statistical decision
theory made differentiable. Wald defined the loss of a decision rule as the
expected payoff under the true distribution, and the optimal procedure is the
minimax rule against the worst admissible distribution. SPO$^+$ and task-based
learning operationalise this by treating the predictor as a parametric family
that minimises an estimate of the decision-induced loss directly, and by training
the family with stochastic gradient descent. The risk bounds of
\citet{LiuGrigas2021} and \citet{HuangGupta2024} turn this into a quantitative
statement about sample complexity. The picture across these strands is that the
classical Granger-Pesaran question, namely which loss to use, now has a
constructive answer at scale, conditional on misspecification and conditional on
the downstream decision being expressible in a differentiable way.

Three obstacles prevent the immediate replacement of standalone forecasting by
decision-focused training. First, TSFM pretraining is action-unconditional, which
means the forecast distribution conditions on observed covariates but not on the
agent's decision. Second, the joint-calibration failure documented by
\citet{PerezDiaz2025} implies that path-dependent decisions, for example
threshold-crossing inventory rules, see less calibrated forecasts than marginal
ones. Third, the Lucas critique applies. A foundation forecaster trained on
observational data cannot distinguish the causal effect of a policy intervention
from the response to a confounding shock, and adding parameters to the forecaster
does not change this.
